% !TITLE 两宋
% !SUMMARY 词与诗并峙的巅峰
% !ORDER 0

\begin{intro}
两宋三百二十年，是中国文学“词”这种文体的黄金时代，也是宋诗自辟蹊径、与唐诗并峙的重要时期。

宋词始于唐而盛于宋。晚唐五代，温庭筠、李煜等人已经把词从民间小曲提升到了文人雅玩的层次，但真正让词成为与诗并驾齐驱的文学体裁的，是宋代。晏殊的闲雅、欧阳修的深婉、柳永的铺陈、苏轼的豪放、周邦彦的精严、辛弃疾的悲壮、李清照的凄切——每一个名字都代表了一种独特的词风。词在宋代完成了从“诗余”到独立文体的身份确认。

苏轼是一个分水岭。在他之前，词多写男女情爱、离愁别绪，所谓“词为艳科”。苏轼以诗为词，“大江东去，浪淘尽，千古风流人物”——词可以写历史、写人生、写哲理，题材和境界一下子打开了。虽然当时有人批评他“不协音律”，但正是这种突破，让词获得了与诗同等的表达力。

宋诗则走了与唐诗不同的路。宋人学唐而变唐，追求理趣、讲究用典、注重句法。“不识庐山真面目，只缘身在此山中”——宋诗的好处在思理的深致，在寻常景物中发掘人生的况味。这种风格，与宋人重理、重学、重内省的文化气质是相契合的。

两宋还是中国历史上文化最为昌明的时代。印刷术的普及、科举制度的完善、文人政治的兴盛，造就了一个全民读书的社会氛围。词在这样的土壤里生根发芽，从宫廷走向市井，从文人走向大众，成为中国文学史上最动人的篇章之一。
\end{intro}
